\documentclass[../report.tex]{subfiles}

\begin{document}

\subsection{Linux - Complete Pentest (Mesa)}

\subsubsection{OSINT}
The following information relating to credentials was found explicitly listed on publicly accessible websites:

\begin{tabular}{l | l | l | l}
\textbf{Employee} & \textbf{Usernames} & \textbf{Passwords} & \textbf{Discoverable by crawling}\\ \hline
Buckerzerg & - & sunshine, tinkerbell & social media\\ 
Groves & sam, groves & machine & mesa homepage\\
Freeman & freeman  & - & mesa homepage\\ 
\end{tabular}

This information being public makes the service vulnerable to username / password enumeration attacks.
\textit{Note:} The username \enquote{marq} could feasibly be guessed.

\textit{Usernames and passwords should not correlate to personal information of preferences. Generating strong passwords and using a password manager is advised. It could also be considered to lower the attack surface by exposing less information publicly and enforcing password guidelines.}

\subsubsection{Active Scan}
Fuzzing directories revealed a publicly accessible backup directory (\textit{/backup}).
This directory contained a database backup which reveled the following information:

\begin{tabular}{l | l | l | r}
\textbf{Webapp-User} & \textbf{Unix-User} & \textbf{Password-Hash} & \textbf{Salary} \\\hline
marq & zucc & 0571749e2ac330a7455809c6b0e7af90 & 10.000.000 \\
freeman & freeman & fddd21b9d7ce17da93c30fa5a653a1df & 250.000 \\
groves & sam & 14754f13e5280c5d49d2ae536c2d57e2 & 133.700 \\
joe & - & c13d23675b7a621212c3a6bb07e0e8df & 80.000 \\
heuzeroth & - & cd1142c62ebbe49a17bc8788a4c83bec & 1 \\[2em]
\end{tabular}

\textit{Autoindexing should be turned off to prevent this from being easily exploited. The directory should not be discoverable and accessible by unauthorized users.}

The collected information was used in the following exploits.

\newpage
\subsubsection{Web Application (Management System) Exploitation}
\paragraph{Guessing / Cracking Credentials}
Analyzing the backup and cracking the hashes revealed the following credentials:

\begin{table}[ht]
\centering
\begin{tabular}{l | l}
\textbf{Usernames} & \textbf{Passwords} \\ \hline
marq & sunshine\\ 
groves & machine \\
freeman  & ?????? \\
joe & hackerman\\
heuzeroth & anwendungssicherheit
\end{tabular}
\end{table}

\paragraph{Brute Forcing Passwords}
The \textit{login.php} page is susceptible to online brute forcing attacks, enabling compromise of passwords with the known account names extracted prior.
\textit{Limiting the amount of repeated login attempts could prevent brute forcing of credentials.}
The credentials for \textit{marq} and \textit{groves} are feasibly guessable without cracking the hashes.

\paragraph{SQL Injection}
Logging into the application is also possible via SQL injection with any known username.\\
This works because user input is not escaped, which means it can alter the SQL statement that is being executed (e.g. \textit{marq'\#}).\\
\textit{Prepared statements properly handle escaping of user input, preventing them from manipulating the SQL query.}

\paragraph{File inclusion and Command Injection}
The PHP management program directly executes user input passed as the \textit{read} URL parameter. This enables users to read files on the system and execute commands. This can be used to create a reverse shell.\\
\textit{Removing the ability to execute commands directly is advised. To retain the ability to read files, validating input and discarding it, if it does not match the intended files, is advised. This way input isn't  executed directly, but only used to determine what should be executed, only if the input is valid.}

\paragraph{SSH Credential Compromise}
Samantha Groves and Gordon Freeman reuse their credentials for their Unix accounts. Marquis Buckerzerg's password can be brute forced by supplying a wordlist extracted from his \textit{x.com} profile, providing root access to the system.
\textit{Passwords should never be reused, and password authentication via SSH should be disabled. Public keys are a more secure authentication method.}

The following Unix account credentials were compromised:
\begin{table}[ht]
\centering
\begin{tabular}{l | l}
\textbf{Username} & \textbf{Password}\\ \hline
zucc & tinkerbell\\
freeman & ??????\\
sam & machine
\end{tabular}
\end{table}

\begin{tabular}{l l}
\textbf{Vector}: & CVSS:4.0/AV:L/AC:L/AT:N/PR:N/UI:N/VC:H/VI:H/VA:H/SC:H/SI:H/SA:H\\
\textbf{CVSS v4.0 Score}: & 9.4 / Critical\\
\end{tabular}

\paragraph{Privilege Escalation}
\textbf{Unsecured Credentials: Shell History (T1552.003)}\\
The user \textit{sam} tried switching to \textit{zucc}, leaving the password readable in plain text in the bash history file.
\textit{Monitoring systems can detect suspicious reads of the bash history and subsequent logins. Preventing credentials from ending up in the bash history is also possible, for example with private history sessions.}

\textbf{Abuse Elevation Control Mechanism (T1548.004)}\\
The user \textit{freeman} has elevated privileges with the program Composer. It can be abused to execute arbitrary code with root privileges.\\
\textit{A regular vulnerability scan could reveal such misconfigurations, enabling them to be corrected.}

\begin{tabular}{l l}
\textbf{Vector}: & CVSS:4.0/AV:L/AC:L/AT:N/PR:L/UI:N/VC:H/VI:H/VA:H/SC:H/SI:H/SA:H\\
\textbf{CVSS v4.0 Score}: & 9.3 / Critical\\
\end{tabular}

\paragraph{phpMyAdmin}
Passwords were reused for this application as well, which should be avoided.
The following\\ phpMyAdming credentials were compromised:

\begin{table}[ht]
\centering
\begin{tabular}{l | l}
\textbf{Username} & \textbf{Password}\\ \hline
marq & sunshine\\
sam & machine
\end{tabular}
\end{table}

The version of phpMyAdmin running on the server has a Remote Code Execution vulnerability, which was abused to create a reverse shell.\\
\textit{Keeping software up to date can mitigate vulnerabilities via security patches.}

\begin{tabular}{l l}
\textbf{Vector}: & CVSS:3.1/AV:N/AC:L/PR:L/UI:N/S:U/C:H/I:H/A:H\\
\textbf{CVSS v3.1 Score}: & 8.8 / High\\
\end{tabular}

\end{document}